% ============================================================
% Section 41: 二维正方晶格的声子色散与空位缺陷散射
% ============================================================
\newpage
\section{固体理论：二维晶格的声子色散与空位缺陷散射}
\label{sec:2d-vacancy-scattering}

\begin{modelbox}[题目：含空位缺陷的二维正方晶格声子]
考虑原子的二维正方形排列，晶格常数为 $a$，原子间有简谐相互作用。当位于 $(x_0,y_0)$ 的原子沿 $x$ 方向位移 $\Delta x$ 时，受到
\begin{itemize}[nosep]
    \item $(x_0\pm a,y_0)$ 两个邻居的恢复力 $-c_1\Delta x$
    \item $(x_0,y_0\pm a)$ 两个邻居的恢复力 $-c_2\Delta x$
\end{itemize}
且 $c_2<c_1$。
\begin{enumerate}
    \item 求沿 $x$ 方向传播的声波色散关系，并画出色散图；
    \item 求此二维晶体中的声速；
    \item 若晶格中存在空位（浓度为 $n$），对平面声波平均恢复力减小量正比于 $n$，且每个空位产生散射幅 $A/\lambda$（$\lambda\gg a$）。求声波的衰减长度作为波长的函数；
    \item 设形成一个空位所需能量为 $\varepsilon$，衰减长度如何依赖于温度？
    \item 在同一假设下，声速如何依赖于温度？
\end{enumerate}
\end{modelbox}

\begin{tipbox}[考点快速分析]
\begin{itemize}
    \item \textbf{单偏振二维格子}：只考虑 $x$ 方向位移，故只有一支纵波色散
    \item \textbf{各向异性}：$x$ 方向邻居贡献 $c_1$，$y$ 方向邻居贡献 $c_2$，方向依赖性进入色散
    \item \textbf{缺陷散射}：空位作为点缺陷，散射幅 $\propto 1/\lambda$，属于 Rayleigh 散射（长波极限）
    \item \textbf{热激活}：空位浓度服从 Boltzmann 分布 $n\propto e^{-\varepsilon/k_BT}$
\end{itemize}
\end{tipbox}

\begin{derivationbox}[标准解题过程]
\textbf{Step 1: 理想晶格的色散关系}

设格点 $(n,m)$ 处原子在 $x$ 方向的位移为 $x_{nm}$。由题给恢复力，运动方程为
\begin{equation}
m\ddot x_{nm} = c_1(x_{n+1,m}+x_{n-1,m}-2x_{nm}) + c_2(x_{n,m+1}+x_{n,m-1}-2x_{nm}).
\label{eq:vacancy-motion}
\end{equation}
设平面波解 $x_{nm}=A\,e^{i(q_1 na+q_2 ma-\omega t)}$，代入并利用 $e^{i\theta}+e^{-i\theta}-2=-4\sin^2(\theta/2)$：
\begin{equation}
-m\omega^2 = -4c_1\sin^2\frac{q_1 a}{2} - 4c_2\sin^2\frac{q_2 a}{2}.
\end{equation}
故普遍色散关系为
\begin{equation}
\boxed{\omega(\vec q) = \sqrt{\frac{4c_1}{m}\sin^2\frac{q_1 a}{2} + \frac{4c_2}{m}\sin^2\frac{q_2 a}{2}}.}
\label{eq:2d-vacancy-dispersion}
\end{equation}

\textbf{Step 2: 沿 $x$ 方向的色散与声速}

令 $q_2=0$，$q_1=q$，定义 $\omega_m=\sqrt{4c_1/m}$，得
\begin{equation}
\boxed{\omega(q) = \omega_m\,\bigl|\sin\tfrac{qa}{2}\bigr|, \qquad 0\le q\le\frac{\pi}{a}.}
\label{eq:x-dispersion-vacancy}
\end{equation}
色散图为从原点出发的正弦半波，边界值 $\omega(\pi/a)=\omega_m$。

声速取群速在长波极限的值。对式 \eqref{eq:x-dispersion-vacancy}，$q\to0$ 时 $\sin(qa/2)\approx qa/2$，故
\begin{equation}
\omega \approx \omega_m\,\frac{qa}{2} = q\,a\sqrt{\frac{c_1}{m}}.
\end{equation}
因此声速
\begin{equation}
\boxed{v_s = a\sqrt{\frac{c_1}{m}}.}
\end{equation}

\textit{注}：若从二维普遍色散 \eqref{eq:2d-vacancy-dispersion} 求群速矢量
\[
\vec v_g = \nabla_{\vec q}\,\omega = \Bigl(\frac{\partial\omega}{\partial q_1},\frac{\partial\omega}{\partial q_2}\Bigr),
\]
在 $q_2=0$、$q_1\to0$ 极限下同样得到 $v_g=a\sqrt{c_1/m}$，与上式一致。

\textbf{Step 3: 空位浓度与衰减长度}

设声波强度为 $I$，传播方向横截面积为 $S$。在距离 $dx$ 内遇到的空位数为 $nS\,dx$。

每个空位的散射幅为 $A/\lambda$，散射截面（正比于散射幅的平方）为
\begin{equation}
\sigma \propto \Bigl(\frac{A}{\lambda}\Bigr)^2 = \frac{A^2}{\lambda^2}.
\end{equation}
强度衰减满足
\begin{equation}
-dI\cdot S \propto I\cdot\sigma\cdot nS\,dx = I\,\frac{A^2}{\lambda^2}\,nS\,dx.
\end{equation}
整理得
\begin{equation}
\frac{dI}{I} = -\frac{c\,nA^2}{\lambda^2}\,dx,
\end{equation}
其中 $c$ 为待定比例常数。积分得 $I(x)=I_0\exp(-x/l)$，衰减长度
\begin{equation}
\boxed{l = \frac{\lambda^2}{c\,nA^2}\propto\frac{\lambda^2}{n}.}
\end{equation}
\textit{物理意义}：波长越长，散射越弱（Rayleigh 散射 $\sigma\propto\lambda^{-2}$），衰减越慢；空位越多，衰减越快。

\textbf{Step 4: 衰减长度与温度的关系}

空位由热激活产生，服从 Boltzmann 分布
\begin{equation}
n \propto e^{-\varepsilon/k_BT}.
\end{equation}
代入衰减长度表达式：
\begin{equation}
\boxed{l \propto \frac{1}{n} \propto e^{\,\varepsilon/k_BT}.}
\end{equation}
温度升高 → 空位浓度指数增长 → 衰减长度指数下降。

\textbf{Step 5: 声速与温度的关系}

存在空位时，有效恢复力系数被削弱。设完整晶体的系数为 $c_1^{(0)},c_2^{(0)}$，则
\begin{equation}
c_1 = c_1^{(0)}(1-Bn), \qquad c_2 = c_2^{(0)}(1-Bn),
\end{equation}
$B$ 为与空位浓度成正比的削弱系数。

声速 $v_s\propto\sqrt{c_1}$，故
\begin{equation}
v_s \propto \sqrt{c_1^{(0)}(1-Bn)} = v_s^{(0)}\sqrt{1-Bn}.
\end{equation}
对稀薄空位 $Bn\ll1$，展开得 $v_s\approx v_s^{(0)}(1-\frac{Bn}{2})$。代入 $n\propto e^{-\varepsilon/k_BT}$：
\begin{equation}
\boxed{v_s(T) \propto \sqrt{1-B'\,e^{-\varepsilon/k_BT}} \approx v_s^{(0)}\Bigl(1-\frac{B'}{2}e^{-\varepsilon/k_BT}\Bigr).}
\end{equation}
温度升高 → 空位增多 → 有效恢复力减弱 → 声速略微下降。低温极限 $T\to0$ 时 $n\to0$，声速恢复完整晶体值 $v_s^{(0)}$。
\end{derivationbox}

\begin{formulabox}[答案汇总]
\begin{center}
\begin{tabular}{cl}
\toprule
问 & 结果 \\
\midrule
(1) 色散关系 & $\displaystyle \omega(q)=\sqrt{\frac{4c_1}{m}}\,\bigl|\sin\tfrac{qa}{2}\bigr|$，$0\le q\le\pi/a$ \\[8pt]
(2) 声速 & $\displaystyle v_s = a\sqrt{\frac{c_1}{m}}$ \\[8pt]
(3) 衰减长度 & $\displaystyle l = \frac{\lambda^2}{c\,nA^2}\propto\frac{\lambda^2}{n}$ \\[8pt]
(4) $l$ 与 $T$ & $\displaystyle l\propto e^{\,\varepsilon/k_BT}$（$T$ 升高，$l$ 指数下降） \\[8pt]
(5) 声速与 $T$ & $\displaystyle v_s(T)\propto\sqrt{1-B'e^{-\varepsilon/k_BT}}$（$T$ 升高，声速略微下降） \\
\bottomrule
\end{tabular}
\end{center}
\end{formulabox}

\begin{insightbox}[物理图像与概念串联]
\textbf{1. 与 sec40 的对比}

本题与 sec40 的数学结构几乎相同（都是二维正方格子、各向异性耦合），但有本质区别：
\begin{itemize}
    \item \textbf{sec40}：考虑 $x$、$y$ 两个偏振方向，得到两支色散（两支声学支）。
    \item \textbf{本题}：只考虑 $x$ 方向位移（平面纵波），得到\textbf{一支}色散。$y$ 方向的邻居只通过 $c_2$ 对 $x$ 振动产生``侧向''恢复力，不引入新的偏振自由度。
\end{itemize}

\textbf{2. Rayleigh 散射的标度律}

点缺陷散射幅 $A/\lambda$ 是 Rayleigh 散射的典型特征：散射体尺度 $\ll\lambda$ 时，散射幅反比于波长。截面 $\sigma\propto\lambda^{-2}$，导致衰减长度 $l\propto\lambda^2$——\textbf{长波声子穿透力更强}。

\textbf{3. 温度依赖的层次}
\begin{itemize}
    \item 空位浓度：$n\propto e^{-\varepsilon/k_BT}$（Arrhenius 型，指数敏感）
    \item 衰减长度：$l\propto 1/n\propto e^{\varepsilon/k_BT}$（同样指数敏感）
    \item 声速：$v_s\propto\sqrt{1-Bn}$（仅通过 $n$ 产生微弱的线性修正）
\end{itemize}
因此，温度对衰减长度的影响远比对声速的影响剧烈。

\textbf{4. 考场注意：区分``群速度''与``声速''}

严格来说，声速是群速度在长波极限的值。若题目直接问``色散关系''，写出 $\omega(q)$ 即可；若问``声速''，必须取 $q\to0$ 极限求 $d\omega/dq$，不能误用相速 $\omega/q$（尽管在本题长波极限二者相等）。
\end{insightbox}

\begin{thinkbox}[考场速记：缺陷散射问题的答题模板]
遇到``点缺陷 + 声波衰减''类问题，按以下步骤展开：
\begin{enumerate}
    \item \textbf{写出色散}：理想晶格 → 平面波代入 → $m\omega^2=\sum_i 4c_i\sin^2(q_i a/2)$
    \item \textbf{求声速}：$q\to0$ 展开 $\omega\approx v_s q$，读出 $v_s$
    \item \textbf{衰减长度}：
    \begin{itemize}
        \item 散射截面 $\sigma\propto |\text{散射幅}|^2$
        \item 衰减系数 $\alpha=n\sigma$
        \item 衰减长度 $l=1/\alpha$
    \end{itemize}
    \item \textbf{温度依赖}：
    \begin{itemize}
        \item 缺陷浓度 $n\propto e^{-\varepsilon/k_BT}$（热激活）
        \item 代入 $l$ 和 $v_s$ 的表达式
    \end{itemize}
\end{enumerate}
\end{thinkbox}
